\section{Maximal Rectangle}
\label{sec:dp-maximal-rectangle}

\problemstatement{Given a binary matrix $\textit{mat}$, find the area
of the largest rectangle (not necessarily square) containing only
$1$s.}

\begin{patternbox}
This closing problem of the entire book splits cleanly into two
independent techniques: a simple per-column \textbf{height DP} (how
many consecutive $1$s sit directly above and including this cell),
followed by treating each row's height array as a \textbf{histogram}
and finding its largest rectangle -- exactly the monotonic-stack
technique already developed for Largest Rectangle in Histogram in the
Stack and Queue chapter.
\end{patternbox}

\subsection*{Deriving the height DP}

Let $\textit{height}[i][j]$ be the number of consecutive $1$s ending at
row $i$, looking upward, in column $j$:
\[
  \textit{height}[i][j] = \begin{cases}
  0 & \textit{mat}[i][j]=0, \\
  \textit{height}[i{-}1][j] + 1 & \textit{mat}[i][j]=1.
  \end{cases}
\]
This is the same ``reset to $0$ on failure, otherwise extend by $1$''
shape seen throughout this book (Longest Common Substring's mismatch
reset, House Robber's running totals) -- here applied independently to
every column.

\subsection*{Reframing each row as a histogram}

After computing $\textit{height}[i][\cdot]$ for a fixed row $i$, that
row \emph{is} a histogram: $\textit{height}[i][j]$ is the bar height at
position $j$. The largest all-ones rectangle with its \textbf{bottom
edge on row $i$} is exactly the largest rectangle in that histogram --
solved in $O(m)$ using the monotonic-stack technique from the Stack and
Queue chapter's Largest Rectangle in Histogram problem. Running this
once per row and tracking the global maximum solves the full problem.

\subsection*{C++ implementation}

\begin{lstlisting}
#include <bits/stdc++.h>
using namespace std;

int largestRectangleInHistogram(vector<int>& heights) {
    int n = heights.size();
    stack<int> st;
    int maxArea = 0;

    for (int i = 0; i <= n; i++) {
        int currentHeight = (i == n) ? 0 : heights[i];
        while (!st.empty() && heights[st.top()] >= currentHeight) {
            int height = heights[st.top()];
            st.pop();
            int width = st.empty() ? i : i - st.top() - 1;
            maxArea = max(maxArea, height * width);
        }
        st.push(i);
    }
    return maxArea;
}

int maximalRectangle(vector<vector<int>>& mat) {
    int n = mat.size(), m = mat[0].size();
    vector<int> height(m, 0);
    int best = 0;

    for (int i = 0; i < n; i++) {
        for (int j = 0; j < m; j++) {
            height[j] = (mat[i][j] == 1) ? height[j] + 1 : 0;
        }
        best = max(best, largestRectangleInHistogram(height));
    }
    return best;
}

int main() {
    vector<vector<int>> mat = {
        {1, 0, 1, 0, 0},
        {1, 0, 1, 1, 1},
        {1, 1, 1, 1, 1},
        {1, 0, 0, 1, 0}
    };
    cout << maximalRectangle(mat) << endl; // 6
    return 0;
}
\end{lstlisting}

\subsection*{Dry run}

\dryrun{Take the $4\times5$ matrix above.}

Row-by-row height arrays: row $0$: $[1,0,1,0,0]$. row $1$:
$[2,0,2,1,1]$. row $2$: $[3,1,3,2,2]$. row $3$: $[4,0,0,3,0]$.

At row $2$, the histogram $[3,1,3,2,2]$ has its largest rectangle as
columns $2$--$4$ at height $2$ (area $6$): bars of height $3,2,2$, the
minimum height $2$ spanning width $3$, giving area $6$. This is the overall maximum: at row $0$ the best is area $1$; at row
$1$, best is $3\times1=3$ (columns $2$--$4$, heights $2,1,1$, limited
by the minimum height $1$); at row $3$, best is $4$ (column $0$ alone,
height $4$). The global maximum across all rows is $6$.

\begin{complexitybox}
\textbf{Time Complexity.} $O(n \cdot m)$ -- $n$ rows, each requiring an
$O(m)$ height update and an $O(m)$ histogram pass (the monotonic stack
visits each bar a bounded number of times).
\textbf{Space Complexity.} $O(m)$ for the \textit{height} array and the
stack.
\end{complexitybox}

\begin{takeawaybox}
This book closes on a problem that is itself a closing argument for
everything in it: recognizing that a 2D problem decomposes into ``a
simple DP recurrence run independently per column'' \emph{composed
with} ``a classical 1D technique from an earlier chapter, run once per
row'' is, in the end, the same skill exercised throughout this entire
book -- noticing which already-solved shape a new problem's structure
matches, rather than deriving a solution from nothing.
\end{takeawaybox}
